\documentclass{article}

\usepackage{arxiv}
\usepackage[utf8]{inputenc}
\usepackage[T1]{fontenc}
\usepackage{hyperref}
\usepackage{url}
\usepackage{booktabs}
\usepackage{amsfonts}
\usepackage{nicefrac}
\usepackage{microtype}
\usepackage{graphicx}
\usepackage{amsmath}
% \graphicspath{ {./images/} {./DAlambert/} {./} }


% override
\renewcommand{\refname}{Referencias}
\renewcommand{\contentsname}{Índice}
\usepackage{fancyhdr}
\fancyhead[R]{\today}
% -----------------------------------

\title{TP 2.2 - GENERADORES DE NÚMEROS PSEUDOALEATORIOS DE DISTINTAS DISTRIBUCIONES DE PROBABILIDAD}
\author{
  Manuel Ferraris \\
  Universidad Tecnológica Nacional - FRRO \\
  Zeballos 1341, S2000, Argentina \\
  \texttt{ferrarismanu@gmail.com} \\
  \And
  Juan Ignacio Fiele \\
  Universidad Tecnológica Nacional - FRRO \\
  Zeballos 1341, S2000, Argentina \\
  \texttt{juanfiele2002@gmail.com} \\
  \And
  María Belén Giannone \\
  Universidad Tecnológica Nacional - FRRO \\
  Zeballos 1341, S2000, Argentina \\
  \texttt{mariabelengiannone010@gmail.com} \\
  \And
  María Lucía Munné \\
  Universidad Tecnológica Nacional - FRRO \\
  Zeballos 1341, S2000, Argentina \\
  \texttt{marilumunne@gmail.com} \\
  \And
  Agustín Stringa \\
  Universidad Tecnológica Nacional - FRRO \\
  Zeballos 1341, S2000, Argentina \\
  \texttt{stringaagustin1@gmail.com} \\
}
\date{\today}
\begin{document}
\maketitle

\keywords{Python \and Numeros Aleatorios \and Distribuciones \and Probabilidad \and Simulación}

\begin{abstract}
\end{abstract}

\section{Introducción}

\section{Desarrollo Teórico}

\subsection{Distribuciones de probabilidad}


\subsection{Pruebas de aleatoriedad}

\section{Resultados}


\section{Discusión y Conclusiones}

\begin{thebibliography}{9}

  \bibitem{enunciado}
  Cátedra de Simulación. (2026). \textit{Enunciado Trabajo Práctico 2.2: GENERADORES DE NÚMEROS PSEUDOALEATORIOS DE DISTINTAS DISTRIBUCIONES DE PROBABILIDAD}. Universidad Tecnológica Nacional - Facultad Regional Rosario (UTN-FRRO).

\end{thebibliography}

\end{document}